% ECONOMICS_PROBLEM: {"id": "O31-207", "title": "Redirecting technological change", "jel_code": "O31", "source_id": "OP-0270"}
% !TeX program = xelatex
\documentclass[11pt,a4paper]{article}
\usepackage[margin=23mm]{geometry}
\usepackage{fontspec}
\setmainfont{Latin Modern Roman}
\usepackage{amsmath,amssymb,mathtools}
\usepackage{xcolor,enumitem,booktabs,longtable,array,xurl,hyperref}
\definecolor{muted}{HTML}{53636C}
\hypersetup{colorlinks=true,urlcolor=blue,linkcolor=blue}
\setlength{\parindent}{0pt}
\setlength{\parskip}{5pt}
\newcommand{\R}{\mathbb R}
\newcommand{\E}{\mathbb E}
\newcommand{\N}{\mathbb N}
\newcommand{\ind}{\mathbf 1}
\newcommand{\Var}{\operatorname{Var}}
\newcommand{\Cov}{\operatorname{Cov}}
\newcommand{\supp}{\operatorname{supp}}
\DeclareMathOperator*{\argmin}{arg\,min}
\DeclareMathOperator*{\argmax}{arg\,max}
\newcommand{\entrymeta}[3]{{\sffamily\small\color{muted}\textbf{\texttt{#1}}\quad|\quad #2\par #3\par}}
\newcommand{\Problemref}[1]{\texttt{#1}}
\newcommand{\JELref}[2]{\texttt{#2} (JEL \texttt{#1})}
\begin{document}
\section*{Redirecting technological change}
\textbf{Problem ID:} \texttt{O31-207}\quad \textbf{JEL:} \texttt{O31}\par
% BEGIN ECONOMICS STATEMENT
\label{problem:OP-0270}
{\sffamily\small\color{muted}Original subject group: Growth and productivity\par}
\label{definitions:problem:30007551}
\textbf{Audit classification:} Explicit published open question. Redirecting innovation while preserving experimentation is expressly posed; the proposed mechanism is editorial and strategic. Verification concerns the cited version, not first priority or proof that this exact editorial specification remains unsolved on October 8, 2026.
\subsection*{Definitions and precise research target}
\paragraph{Editorial mathematical specification.} Fix two research directions \(j=1,2\), such as ordinary private productivity and a declared socially valuable health or environmental technology. Firms privately observe project quality \(q\), choose direction and research effort \(r\), and obtain innovation arrival rate \(\lambda_j(r,q)\). Private rewards \(v_j(q)\) differ from social rewards \(V_j(q)\), including measured spillovers and external damages. A policymaker observes imperfect signals \(s\), can offer subsidies \(t_j(s,r)\), and faces a stated budget and an explicit probability of political capture or biased project selection.
Determine a feasible subsidy or procurement rule maximizing expected discounted social surplus net of research resources, fiscal costs, and capture losses. Require firms’ participation and effort/direction incentive compatibility for every type in the declared quality support. Preserve a specified minimum experimentation probability for projects with uncertain social value, or incorporate the information value of experimentation directly in the objective.
Identify the private-response, spillover, learning, and capture parameters needed for policy comparison, and distinguish uncertainty about project returns from discretionary favoritism. Evaluate rules against untargeted support and no-support benchmarks under equal budgets. The target is direction-specific innovation policy with explicit informational constraints; it does not assert that a policymaker knows socially optimal technology ex ante or that one favored sector should receive support universally.

Define an innovation within a finite interval \(\Delta>0\) by probability \(p_j(r,q)=1-e^{-\lambda_j(r,q)\Delta}\), with \(\lambda_j\ge0\); this prevents confusing an arrival intensity with a probability. Let a firm maximize \(p_j(r,q)v_j(q)-c_j(r,q)+t_j(s,r)\) over direction and effort after observing quality. Incentive compatibility means the prescribed choice solves this maximization for every quality and observed signal, and report-based schemes require additional truthfulness. For a static benchmark social surplus is \(E[p_j(r,q)V_j(q)-c_j(r,q)-C^F(t)-L^{\rm cap}(t)]\), with subsidies counted as transfers except for real financing and capture losses. Experimentation is a positive allocation to a defined uncertain-project set, e.g. \(P(j\in U,r\ge r_0)\ge\epsilon\); in a dynamic learning model its option value instead enters continuation welfare. Direction/private-information incentives overlap game theory.
\subsection*{Variants and assumption changes}
The following are \emph{editorial variants}, not additional published open conjectures.
\begin{enumerate}
\item \emph{Public research quality.} Quality and social value are observable, there is no capture, and effort is verifiable. Derive a direction-specific subsidy benchmark for the private/social return wedge, with a binding common budget.
\item \emph{Learning and capture.} Quality is initially unknown to the policymaker, pilot outcomes update a posterior, and biased allocation follows a specified probability kernel. Optimize adaptive subsidies subject to experimentation and fiscal limits and compare with untargeted support.
\end{enumerate}
\subsection*{References and source audit}
\begin{itemize}
\item Ufuk Akcigit. \emph{Creative destruction in economic growth}. 2026. Locator: Section7, Open questions, item7. Primary text checked. \url{https://onlinelibrary.wiley.com/doi/10.1111/sjoe.70037}.
\end{itemize}
Redirecting innovation while preserving experimentation is expressly posed; the proposed mechanism is editorial and strategic. This formulation contains strategic incentives and therefore overlaps game theory.
\begin{enumerate}\raggedright
\item\hypertarget{definitions:source:30007551:1}{} Ufuk Akcigit. 2026. \emph{Creative destruction in economic growth}. Section7, Open questions, item7.
\url{https://onlinelibrary.wiley.com/doi/10.1111/sjoe.70037}.
DOI: \url{https://doi.org/10.1111/sjoe.70037}.
\end{enumerate}

\subsection*{Common conventions: Source mathematical conventions}

These conventions apply to each entry unless explicitly replaced.

\textbf{Models and identification.} A primitive model \(m\) specifies all probability laws, feasible choices, information, equilibrium and measurement rules used by its target. The admissible class \(\mathcal M\) is fixed independently of outcomes. If \(P_O(m)\) is its observable probability law and \(\tau(m)\) the target, its identified set at observable law \(P\) is
\[
 \mathcal I_\tau(P)=\{\tau(m):m\in\mathcal M,\ P_O(m)=P\}.
\]
Point identification means that this set is a singleton. An empty set signals incompatibility of data and assumptions. Coordinate bounds need not describe the joint set. Bounds are sharp only if every retained target is attainable or arbitrarily approachable by compatible models. Finite bounds require explicit restrictions on unobserved quantities; unrestricted confounding, hidden wealth, extrapolation or arbitrary equilibrium selection can yield uninformative sets.

\textbf{Causal interventions.} A potential outcome \(Y_i(p)\) is the outcome under a complete policy regime \(p\), including timing, financing, information and market spillovers. The target population and units are fixed before treatment. With interference, use the complete assignment vector or a justified exposure mapping; individual no-interference notation alone cannot identify an economy-wide intervention. A difference of mean potential outcomes does not require the cross-world joint distribution. Conditioning on post-treatment migration, survival, enrollment or employment changes the estimand and needs separate justification.

\textbf{Statistical statements.} An estimator, confidence set, forecast or test specifies a sampling law, model class and loss/coverage criterion. Uniform \((1-\alpha)\)-coverage means \(\inf_{m\in\mathcal M}\Pr_m\{\tau(m)\in C_n\}\geq1-\alpha\), or an explicitly asymptotic version. The whole parameter space trivially has coverage, so informativeness requires a stated width, risk or power objective. A forecasting improvement is relative to a named benchmark, information set and evaluation population.

\textbf{Equilibrium and optimization.} An equilibrium comprises feasible plans, optimality, beliefs where needed, budget constraints and all market/resource clearing. The primitives or a stated benchmark define each correspondence \(\mathcal E(p;m)\). Nonemptiness is assumed or proved, never inferred just from compact policy sets. A selected-equilibrium target uses a declared selection rule; a robust target quantifies over all permitted equilibria. Use suprema or infima unless attainment is established. An \(\varepsilon\)-optimal policy is feasible and within \(\varepsilon\) of the finite optimum value. Report an empty feasible set separately.

\textbf{Welfare and accounting.} Normative utility, interpersonal normalization, social weights, numeraire, discounting and the use of public receipts are primitives. Behavioral utility need not coincide with the welfare criterion. Household welfare includes ownership income and rents once; adding profits again to compensating variation generally double-counts them. A monetary equivalent is defined by an explicit transfer and price convention, with monotonicity and range conditions. Average effects, mechanism contributions, transfers and welfare losses are different targets. Infinite sums require convergence and dynamic feasible plans require terminal or no-Ponzi conditions.

\textbf{Source versions.} A working-paper date, revision date and journal publication year are distinguished. A DOI for a journal version is not automatically the DOI of an earlier manuscript. Locators refer to the stated version and printed pages where specified. A metadata-only check does not verify a claim about a full-text conclusion. Every entry's source subsection narrows the provenance claim accordingly.
% END ECONOMICS STATEMENT
\end{document}
